\chapter{The I-EIP monitor study}\label{ch:ieip}
% src: docs/PREREG_IEIP_TWIN.md
% src: C:\source\erisml-lib\docs\I-EIP_Monitor_Whitepaper.md
% src: twin/ieip/analysis.py, twin/ieip/gpu_capture.py, twin/ieip/nrp.py

The containment of Chapter~\ref{ch:lean} keeps the robot's language models away from authority.
It does not make their readings of the world correct.
The robot's classifier can still report a fall that did not happen or miss one that did.
Such an error cannot by itself grant authority, but it can make the robot check in when it
need not, or fail to ask when it should.
This chapter describes a preregistered study of one way to catch such errors.
The study asks whether a check on the classifier's hidden states predicts that its output is
wrong.
The study was registered on 2 October 2026, before any activation was captured, any replay
classification was produced or the ground-truth labeller was written.
% src: docs/PREREG_IEIP_TWIN.md:3-6
At the time of writing its results are not available.
Every outcome below is a slot to be filled from the graded record.

\section{The question is whether internal inconsistency predicts error}

The robot's classifier turns perception facts into scene events.
Its weakness is a plausible output produced for the wrong internal reason.
% src: docs/PREREG_IEIP_TWIN.md:10-11
The Internal Epistemic Invariance Principle (I-EIP) gives a criterion for such failures
\cite{bond2026ieip}.
Let $x$ be an input, $g$ a transform that rewrites $x$ without changing its meaning, and
$h_\ell(x)$ the model's hidden state at layer $\ell$.
The criterion says the hidden state should change by a fixed linear map $\rho_\ell(g)$ that
depends on the transform and the layer but not on the input.
\begin{equation}\label{eq:ieip}
  h_\ell(g\,x) \;\approx\; \rho_\ell(g)\, h_\ell(x).
\end{equation}
% src: C:\source\erisml-lib\docs\I-EIP_Monitor_Whitepaper.md:36-42
% src: docs/PREREG_IEIP_TWIN.md:11-14
The criterion comes from section 16 of the GUASS-SAI report \cite{bond2025guass}.
The I-EIP Monitor whitepaper extends it from monitoring to gating.
% src: C:\source\erisml-lib\docs\I-EIP_Monitor_Whitepaper.md:16
The whitepaper argues that the internal form is stronger than a check on outputs.
A model whose outputs are invariant under a rewrite but whose hidden states are not may be
reaching the same answer by a different computation.
% src: C:\source\erisml-lib\docs\I-EIP_Monitor_Whitepaper.md:44-50
The whitepaper also lists what the monitor cannot detect.
A model whose hidden states stay consistent while its goal is wrong passes the check.
So do transforms outside the declared set.
% src: C:\source\erisml-lib\docs\I-EIP_Monitor_Whitepaper.md:332-343

The study tests two claims.
The first is that a failure of equation~\eqref{eq:ieip}, measured on the classifier's own hidden
states, predicts that its classification is wrong.
The second is that it says more than a check of the outputs alone.
% src: docs/PREREG_IEIP_TWIN.md:14-16

\paragraph{The prior favours failure.}
The preregistration states its prior before any data.
A consumer-relative variant of this monitor, called CR-IEIP, failed all three of its sealed
families.
No ordering variable predicted when the monitor helped.
A failure here is therefore the expected outcome, and it will be reported as plainly as a pass.
% src: docs/PREREG_IEIP_TWIN.md:18-21

\section{The classifier under test is the robot's on-board model}

The classifier under test is Qwen2.5-7B-Instruct \cite{qwen2024qwen25} in bfloat16.
It runs as the robot's scene classifier, with the \erisml{} compiler's
\code{ObservationClassifier} prompt and the event vocabulary of the scene file.
Decoding is greedy, with at most 512 new tokens.
% src: docs/PREREG_IEIP_TWIN.md:25-27
% src: twin/ieip/analysis.py:41,57
This is the on-robot classifier the robot reasons with when communications are down.
The cloud classifier exposes no activations and is not tested.
% src: docs/PREREG_IEIP_TWIN.md:27-29
The run uses the model revision held in Atlas's Hugging Face cache, \code{a09a354}.
% src: docs/PREREG_IEIP_TWIN.md:205
% src: twin/ieip/nrp.py:82-83

The inputs are the perception facts of every robot decision cycle recorded in the brain's
hash chain.
The facts are replayed and nothing is re-simulated.
% src: docs/PREREG_IEIP_TWIN.md:30-32
Amendment A5 fixes the data as three development runs, \code{dev7}, \code{dev8r} and
\code{dev9r}.
Two earlier runs, \code{dev8} and \code{dev9}, were void.
The brain on Atlas had run a stale library whose DEME gate raised on every request, so the
robot never acted and those runs hold no decision cycles.
% src: docs/PREREG_IEIP_TWIN.md:181-187
The three results files are committed with their SHA-256 hashes, and the replay reads them only
from that commit.
% src: twin/ieip/inputs/SHA256SUMS
% src: twin/ieip/nrp.py:101
The prompt for every cycle is built by the current isolated classifier.
For \code{dev8r} and \code{dev9r} the replayed prompt is byte-identical to the live one.
For \code{dev7}, which ran before the scene agent was isolated, it is not.
Replay is also history-free, since a cycle's recent events are not logged with it.
% src: docs/PREREG_IEIP_TWIN.md:145-148
% src: docs/PREREG_IEIP_TWIN.md:188-192
The number of decision cycles in the three runs is \RESULT{cycles per run and in total}.

\section{Five rewrites, four layers and one score per cycle}

\subsection{The transforms}

Each cycle's facts $x$ are rewritten by five deterministic transforms, seeded by the cycle's
hash.
% src: docs/PREREG_IEIP_TWIN.md:40-41
\begin{description}
  \item[$g_0$, the null transform.] The same JSON with different whitespace and indentation.
  It changes tokens but not meaning and sets the noise floor.
  \item[$g_1$.] Every object's keys in a different order.
  \item[$g_2$.] Distances in centimetres instead of metres.
  \item[$g_3$.] The clock time in twelve-hour form.
  \item[$g_4$.] Fixed synonyms for the simulator's activity and behaviour labels, from a table
  committed before any replay.
  A label not in the table is left as it is.
\end{description}
% src: docs/PREREG_IEIP_TWIN.md:43-50
% src: twin/ieip/analysis.py:60-77 (SYNONYMS)
% src: twin/ieip/analysis.py:91-140 (g0 to g4)
The synonym table maps, for example, \code{sleep} to \code{asleep} and \code{reading} to
\code{reading a book}.
% src: twin/ieip/analysis.py:61,70
Each cycle thus yields six prompts, the original and its five rewrites.
% src: twin/ieip/analysis.py:143-150 (transformed returns x and the five rewrites)
Figure~\ref{fig:ieip-transforms} shows the five rewrites.

\begin{figure}[t]
  \centering
  % Chapter 12. The five meaning-preserving rewrites of a cycle's perception facts.
% src: docs/PREREG_IEIP_TWIN.md:40-50
% src: twin/ieip/analysis.py:60-77 (SYNONYMS), 91-150 (g0 to g4, seeded per cycle)
\begin{tikzpicture}[house,
  p/.style={minimum width=29mm, minimum height=40mm, text width=27mm, anchor=north},
  d/.style={detail, text width=25mm}]
\foreach \i/\col/\x/\lab in {0/Grey/0/g_0, 1/Purple/31/g_1, 2/Blue/62/g_2, 3/Teal/93/g_3, 4/Green/124/g_4}{
  \node[panel=\col, p] (P\i) at (\x mm,0) {};
  \node[step=\col] at ([yshift=-5mm]P\i.north) {\i};
}
\node[stepname=Grey, below=10mm of P0.north] (n0) {$g_0$ null};
\node[desc, below=0mm of n0, text width=27mm] (e0) {whitespace only};
\node[d, below=1.5mm of e0] {indentation\\1, 2 or 4\\sets the noise\\floor};

\node[stepname=Purple, below=10mm of P1.north] (n1) {$g_1$ key order};
\node[desc, below=0mm of n1, text width=27mm] (e1) {every object};
\node[d, below=1.5mm of e1] {keys shuffled\\values unchanged};

\node[stepname=Blue, below=10mm of P2.north] (n2) {$g_2$ units};
\node[desc, below=0mm of n2, text width=27mm] (e2) {metres to cm};
\node[d, below=1.5mm of e2] {key suffix\\\code{\_m} to \code{\_cm}\\value times 100};

\node[stepname=Teal, below=10mm of P3.north] (n3) {$g_3$ clock};
\node[desc, below=0mm of n3, text width=27mm] (e3) {twelve-hour form};
\node[d, below=1.5mm of e3] {\code{11:53}\\becomes\\\code{11:53 AM}};

\node[stepname=Green, below=10mm of P4.north] (n4) {$g_4$ synonyms};
\node[desc, below=0mm of n4, text width=27mm] (e4) {activity, behaviour};
\node[d, below=1.5mm of e4] {\code{sleep} to \code{asleep}\\\code{reading} to\\\code{reading a book}};

% the original facts feed all five
\node[panel=Amber, anchor=south, minimum width=60mm] (X) at ([yshift=8mm]P2.north) {\textbf{facts $x$ of one decision cycle}\quad seed from the cycle's hash};
\foreach \i in {0,...,4}{ \draw[flow] (X.south) -- (P\i.north); }
\end{tikzpicture}

  \caption{The five rewrites applied to the perception facts of each decision cycle. Each keeps
  the meaning of the facts and changes their tokens. The null rewrite changes only whitespace and
  sets the noise floor.}
  \label{fig:ieip-transforms}
\end{figure}

\subsection{The representation and the error}

The representation is the hidden state of the last prompt token at four layers.
They sit at 25, 50, 75 and 90 percent of the model's depth, rounded down.
% src: docs/PREREG_IEIP_TWIN.md:54-55
% src: twin/ieip/gpu_capture.py:37-39
Amendment A2 projects each layer's states onto the top 32 principal components of the
calibration half's untransformed states before anything else is fitted.
The model's hidden states have 3584 dimensions and the calibration half has a few hundred
cycles, so a full map would be badly underdetermined.
% src: docs/PREREG_IEIP_TWIN.md:132-136
% src: twin/ieip/analysis.py:44 (PCA_K = 32)
On the projected states the map $\rho_\ell(g)$ is fitted by regularized least squares on the
calibration half and then frozen.
% src: docs/PREREG_IEIP_TWIN.md:56-57
% src: C:\source\erisml-lib\src\erisml\ieip\rho.py:73-95 (estimate_rho)
The whitepaper writes the estimate as follows, with $X$ and $Y$ the stacked states of the
original and rewritten inputs and $\lambda > 0$ a ridge coefficient.
\begin{equation}
  \hat\rho_\ell(g) \;=\; Y X^{\top} \bigl( X X^{\top} + \lambda I \bigr)^{-1}.
\end{equation}
% src: C:\source\erisml-lib\docs\I-EIP_Monitor_Whitepaper.md:187-189
The error of cycle $x$ under transform $g$ at layer $\ell$ is the relative residual.
\begin{equation}
  e(x, g, \ell) \;=\; \frac{\bigl\lVert h_\ell(g\,x) - \hat\rho_\ell(g)\, h_\ell(x) \bigr\rVert}
  {\lVert h_\ell(x) \rVert}.
\end{equation}
% src: docs/PREREG_IEIP_TWIN.md:58
% src: docs/PREREG_IEIP_TWIN.md:137-139
% src: twin/ieip/analysis.py:367-369

\subsection{The score and the flag}

Errors are measured in units of their own noise.
For each transform $g$ and layer $\ell$, let $m_{g,\ell}$ be the median of $e(\cdot, g, \ell)$
on the calibration half and $s_{g,\ell}$ its median absolute deviation.
The score of a cycle is its largest standardized error over the four real rewrites and the
four layers.
\begin{equation}
  z(x) \;=\; \max_{g \in \{g_1, \dots, g_4\}} \; \max_{\ell} \;
  \frac{e(x, g, \ell) - m_{g,\ell}}{s_{g,\ell}}.
\end{equation}
A cycle is flagged when $z(x) > 3$.
% src: docs/PREREG_IEIP_TWIN.md:156-162 (amendment A3)
% src: twin/ieip/analysis.py:45 (FLAG_Z = 3.0)
% src: twin/ieip/analysis.py:371-384
The registration first measured every transform against the null transform's noise.
A dry run on synthetic states showed that any real rewrite leaves a residual above the null
floor, so every cycle flagged.
Amendment A3 replaced it with each transform's own noise.
The null transform's score is still reported per cycle as a sanity check.
% src: docs/PREREG_IEIP_TWIN.md:154-162
Figure~\ref{fig:ieip-scoring} shows the path from hidden states to the registered tests.

\begin{figure}[t]
  \centering
  % Chapter 12. From hidden states to the registered tests.
% src: docs/PREREG_IEIP_TWIN.md:54-96, 132-139, 149-164, 168-176
% src: twin/ieip/analysis.py:41-56, 331-433
\begin{tikzpicture}[house,
  p/.style={minimum width=23mm, minimum height=44mm, text width=21mm, anchor=north},
  d/.style={detail, text width=19mm}]
\foreach \i/\col/\x in {1/Purple/0, 2/Blue/25.5, 3/Teal/51, 4/Green/76.5, 5/Amber/102, 6/Orange/127.5}{
  \node[panel=\col, p] (P\i) at (\x mm,0) {};
  \node[step=\col] at ([yshift=-5mm]P\i.north) {\i};
}
\node[stepname=Purple, below=10mm of P1.north] (n1) {States};
\node[desc, below=0mm of n1, text width=21mm] (e1) {$x$ and 5 rewrites};
\node[d, below=1.5mm of e1] {last prompt\\token, layers at\\25, 50, 75 and\\90 percent};

\node[stepname=Blue, below=10mm of P2.north] (n2) {Project};
\node[desc, below=0mm of n2, text width=21mm] (e2) {per layer};
\node[d, below=1.5mm of e2] {top 32\\components of\\calibration $x$\\(A2)};

\node[stepname=Teal, below=10mm of P3.north] (n3) {Fit $\hat\rho$};
\node[desc, below=0mm of n3, text width=21mm] (e3) {per rewrite, layer};
\node[d, below=1.5mm of e3] {ridge least\\squares on the\\calibration half\\then frozen};

\node[stepname=Green, below=10mm of P4.north] (n4) {Noise};
\node[desc, below=0mm of n4, text width=21mm] (e4) {per rewrite, layer};
\node[d, below=1.5mm of e4] {median and\\median absolute\\deviation of $e$\\(A3)};

\node[stepname=Amber, below=10mm of P5.north] (n5) {Score};
\node[desc, below=0mm of n5, text width=21mm] (e5) {one $z$ per cycle};
\node[d, below=1.5mm of e5] {largest over\\$g_1$ to $g_4$ and\\layers\\flag if $z > 3$};

\node[stepname=Orange, below=10mm of P6.north] (n6) {Test};
\node[desc, below=0mm of n6, text width=21mm] (e6) {H1 and H2};
\node[d, below=1.5mm of e6] {cluster bootstrap\\over scenarios\\10{,}000 draws\\5th percentile};

\foreach \a/\b in {1/2, 2/3, 3/4, 4/5, 5/6}{
  \draw[flow] ([yshift=-22mm]P\a.north east) -- ([yshift=-22mm]P\b.north west);
}
% the split
\node[panel=Teal, anchor=north, minimum width=74mm, minimum height=7mm] (cal) at ([yshift=-6mm]P3.south) {calibration half\quad odd-numbered scenarios};
\node[panel=Orange, anchor=north, minimum width=48mm, minimum height=7mm] (tst) at ([yshift=-6mm, xshift=12.75mm]P5.south) {test half\quad even-numbered scenarios};
\draw[flow] (cal.north -| P2.south) -- (P2.south);
\draw[flow] (cal.north -| P3.south) -- (P3.south);
\draw[flow] (cal.north -| P4.south) -- (P4.south);
\draw[flow] (tst.north -| P5.south) -- (P5.south);
\draw[flow] (tst.north -| P6.south) -- (P6.south);
\end{tikzpicture}

  \caption{From hidden states to the registered tests. The projection, the maps and the noise
  units are fitted on the calibration half only and then frozen. The test half is scored with
  them and is the only data the hypotheses see.}
  \label{fig:ieip-scoring}
\end{figure}

\subsection{Output invariance and ground truth}

A cycle is output-invariant when the classifier's event set is the same on $x$ and on all five
rewrites.
% src: docs/PREREG_IEIP_TWIN.md:63-64
% src: twin/ieip/analysis.py:388
Ground truth comes from a labelling function on the world state the cycle's facts record.
It labels the event types the world settles unambiguously.
They are falls, harm and its severity band, attacks by animals and by people, smoke, unexpected
visitors, media content and unresponsiveness.
% src: docs/PREREG_IEIP_TWIN.md:65-69
Amendment A1 dropped the two check-in types, which are system events the classifier can never
emit.
It also allowed a label to be unsettled, in which case the cycle is excluded for that type.
% src: docs/PREREG_IEIP_TWIN.md:116-127
A classification is an error when the classifier's events, restricted to the labelled types,
contain a false event or miss a true one.
% src: docs/PREREG_IEIP_TWIN.md:70-71

\section{Three hypotheses on a split fixed by scenario}

\subsection{The split}

Cycles are split by scenario, so no scenario contributes to both halves.
Odd-numbered development scenarios calibrate and even-numbered ones test.
% src: docs/PREREG_IEIP_TWIN.md:75-77
% src: twin/ieip/analysis.py:50-52 (in_calibration)
The registration first sent every later scenario to the test half.
Amendment A4 continued the parity rule instead.
The first run produced 26 decision cycles over 21 scenarios, too few for the test half to reach
its floor.
Eight longer scenarios were added, and under the first rule the calibration half would have had
fewer cycles than the 32 components of the projection.
The assignment is fixed by scenario number, before any run of the new scenarios.
% src: docs/PREREG_IEIP_TWIN.md:166-176

\subsection{The hypotheses}

All tests are on the test half and one-sided.
Each uses a cluster bootstrap over scenarios, the unit the covariates live on, with 10{,}000
resamples and seed 20261002.
% src: docs/PREREG_IEIP_TWIN.md:81-82
% src: twin/ieip/analysis.py:46-47
The one-sided 95 percent lower bound is the 5th percentile of the bootstrap distribution.
% src: docs/PREREG_IEIP_TWIN.md:149-150
% src: twin/ieip/analysis.py:390-433 (_bootstrap_diff)
Let $\Delta$ be the error rate among flagged cycles minus the error rate among unflagged cycles.
\begin{description}
  \item[H1, detection.] On the test half $\Delta > 0$.
  It passes if the lower bound of $\Delta$ is above zero.
  \item[H2, beyond the outputs.] The same difference and rule, among output-invariant cycles
  only.
  This is the claim that the internal check sees what an output check does not.
  \item[H3, gating end to end.] Run only if H1 passes.
  The on-board classifier runs live with gating, so a flagged cycle steps no classified events
  and the robot checks in instead.
  A sealed, unread scenario set is run twice, gating off and on, with everything else frozen at
  one commit.
  It passes if, with gating on, false clears are no more frequent, over-restriction is no more
  frequent and class-correct is at most one scenario lower.
\end{description}
% src: docs/PREREG_IEIP_TWIN.md:84-92

\subsection{The inconclusive rule}

A hypothesis is graded inconclusive, not pass or fail, when its test cycles number fewer than
50, its errors fewer than 5 or its flagged cycles fewer than 5.
% src: docs/PREREG_IEIP_TWIN.md:94-96
% src: twin/ieip/analysis.py:55 (MIN_CYCLES, MIN_ERRORS, MIN_FLAGGED)
Amendment A3 adds that a test half with no flagged or no unflagged cycles leaves $\Delta$
undefined, which is also inconclusive.
% src: docs/PREREG_IEIP_TWIN.md:163-164
% src: twin/ieip/analysis.py:418-425
Whatever the outcome, the report gives the base error rate, the flag rate, the count in each
cell, the interval for each hypothesis, the per-transform and per-layer error distributions and
every flagged cycle with its scenario.
It also reports the whitepaper's named failure modes, \code{group\_symmetry\_break} and
\code{layerwise\_drift}, where they apply.
% src: docs/PREREG_IEIP_TWIN.md:100-102

\section{The capture ran on NRP A10 GPUs with weights on parallel volumes}

The analysis has three stages.
Prompts are built on a CPU.
GPU shards capture hidden states and generate classifications.
The shards are then merged and graded on a CPU.
% src: twin/ieip/analysis.py:2-17
% src: docs/PREREG_IEIP_TWIN.md:193-202

The prompts are built twice from the same commit, once on the National Research Platform's
Nautilus cluster and once on Atlas.
They must match byte for byte, by SHA-256, before any GPU job is sent.
% src: docs/PREREG_IEIP_TWIN.md:194-195
Each GPU shard loads the model once.
The compiler's \code{HuggingFaceActivationSource} wraps the loaded model's base, so the hidden
states are captured with the same hooks and forward pass without a second copy of the weights.
The classification is then generated greedily from the same text in a second pass.
% src: twin/ieip/gpu_capture.py:5-17
% src: docs/PREREG_IEIP_TWIN.md:196-200
The capture reads up to 8192 prompt tokens, since the library's default of 512 would cut the
facts out of the prompt.
% src: twin/ieip/analysis.py:56
% src: docs/PREREG_IEIP_TWIN.md:140-144
The states are stored as their exact bfloat16 bit patterns, and this exactness is checked.
A shard's outputs are written under temporary names and renamed last, so a shard is either
complete or absent.
% src: twin/ieip/gpu_capture.py:11-16
% src: docs/PREREG_IEIP_TWIN.md:204

The GPUs are NVIDIA A10s in the UNL zone of the cluster.
% src: twin/ieip/nrp.py:53-54
The weights first sat on a shared CephFS volume, which pods read at 14.9 MB/s on one stream and
38.3 MB/s on three.
Each A10 would have idled for about seven minutes while loading, which the cluster's rules
forbid.
% src: docs/PREREG_IEIP_TWIN.md:211-214
% src: twin/ieip/nrp.py:13-16
Four linstor block volumes read in parallel gave 578.5 MB/s, against 149.2 MB/s for one.
Each GPU pod therefore mounts its own set of four volumes, holding one weight file each, and
loads the four files in parallel.
Volume 0 also holds the environment, the code, the prompts and that pod's outputs.
% src: docs/PREREG_IEIP_TWIN.md:215-219
% src: twin/ieip/nrp.py:64-66 (LIN_CLASS, SETS, WEIGHT_FILES)
The weights are the same files from the same revision, verified file by file against the
hashes of Atlas's cache.
% src: docs/PREREG_IEIP_TWIN.md:205-206
Amendment A6 records that this changes only where bytes are read from.
% src: docs/PREREG_IEIP_TWIN.md:220-221
Figure~\ref{fig:ieip-nrp} shows the steps of the capture and the volumes of one GPU pod.

\begin{figure}[t]
  \centering
  % Chapter 12. The capture on NRP Nautilus and the layout of a GPU pod's volumes.
% src: twin/ieip/nrp.py:1-33 (steps, layout, rules), 53-66, 101, 107-113
% src: docs/PREREG_IEIP_TWIN.md:193-223 (A5, A6)
% src: twin/ieip/gpu_capture.py:1-17
\begin{tikzpicture}[house,
  p/.style={minimum width=23mm, minimum height=50mm, text width=21mm, anchor=north},
  d/.style={detail, text width=19mm},
  vol/.style={detail, text width=26mm, minimum height=14mm, draw=hsGreenLine, line width=0.6pt}]
\foreach \i/\col/\x in {1/Purple/0, 2/Grey/25.5, 3/Blue/51, 4/Green/76.5, 5/Teal/102, 6/Orange/127.5}{
  \node[panel=\col, p] (P\i) at (\x mm,0) {};
  \node[step=\col] at ([yshift=-5mm]P\i.north) {\i};
}
\node[stepname=Purple, below=10mm of P1.north] (n1) {Inputs};
\node[desc, below=0mm of n1, text width=21mm] (e1) {three runs};
\node[d, below=1.5mm of e1] {dev7, dev8r\\dev9r\\committed with\\SHA-256};

\node[stepname=Grey, below=10mm of P2.north] (n2) {Stage};
\node[desc, below=0mm of n2, text width=21mm] (e2) {CPU jobs};
\node[d, below=1.5mm of e2] {env, code and\\weights verified\\file by file\\onto each set};

\node[stepname=Blue, below=10mm of P3.north] (n3) {Prepare};
\node[desc, below=0mm of n3, text width=21mm] (e3) {CPU, twice};
\node[d, below=1.5mm of e3] {on NRP and on\\Atlas, prompts\\must match by\\SHA-256};

\node[stepname=Green, below=10mm of P4.north] (n4) {Capture};
\node[desc, below=0mm of n4, text width=21mm] (e4) {A10 GPUs at UNL};
\node[d, below=1.5mm of e4] {pilot, then\\shards, at most\\four pods\\states, then reply};

\node[stepname=Teal, below=10mm of P5.north] (n5) {Collect};
\node[desc, below=0mm of n5, text width=21mm] (e5) {fetch to Atlas};
\node[d, below=1.5mm of e5] {every prompt\\exactly once\\exact bf16\\bit patterns};

\node[stepname=Orange, below=10mm of P6.north] (n6) {Grade};
\node[desc, below=0mm of n6, text width=21mm] (e6) {committed script};
\node[d, below=1.5mm of e6] {H1, H2\\written to\\\code{RESULTS\_}\\\code{IEIP\_TWIN.md}};

\foreach \a/\b in {1/2, 2/3, 3/4, 4/5, 5/6}{
  \draw[flow] ([yshift=-21mm]P\a.north east) -- ([yshift=-21mm]P\b.north west);
}

% one GPU pod and its own set of four linstor volumes
\node[vol] (v0) at (16mm,-74mm) {volume 0\\weight file 1\\env, code, prompts\\outputs};
\node[vol, right=4mm of v0] (v1) {volume 1\\weight file 2};
\node[vol, right=4mm of v1] (v2) {volume 2\\weight file 3};
\node[vol, right=4mm of v2] (v3) {volume 3\\weight file 4};
\node[panel=Green, text width=40mm] (pod) at ($(v1.south)!0.5!(v2.south)+(0,-14mm)$) {\textbf{GPU pod}\quad one A10\\loads the four files in parallel};
\foreach \v in {v0,v1,v2,v3}{ \draw[flow] (\v.south) -- (pod.north); }
\begin{scope}[on background layer]
  \node[panel=Grey, fit=(v0)(v3)(pod), inner sep=3mm,
        label={[stepname=Grey, anchor=south west]north west:One set of four linstor volumes per pod (A6)}] (set) {};
\end{scope}
\end{tikzpicture}

  \caption{The capture on the Nautilus cluster. CPU jobs stage the weights and build the
  prompts, A10 pods capture states and replies, and the outputs return to Atlas for grading.
  Each pod reads its weights from its own four block volumes in parallel, so the GPU does not
  wait idle on a slow shared volume.}
  \label{fig:ieip-nrp}
\end{figure}

The submitter scores each job against the cluster's rules in code rather than by recall.
CPU jobs use one core and 2 GiB.
A GPU job installs and downloads nothing, since CPU jobs stage everything first.
A pilot shard is sized by a stated model, and later shards by the pilot's measured usage.
Requests equal limits, every command is bounded by a timeout and at most four GPU pods run at
once.
% src: twin/ieip/nrp.py:21-33
% src: twin/ieip/nrp.py:107-113
The pilot's measured load time is \RESULT{load time from the first GPU pod's log}.
% src: docs/PREREG_IEIP_TWIN.md:221-223

\section{Results}

No result is available at the time of writing.
The committed script will write the graded record to the repository file
\code{docs/RESULTS\_IEIP\_TWIN.md}.
% src: docs/PREREG_IEIP_TWIN.md:111
The slots below follow the list the preregistration requires.
% src: docs/PREREG_IEIP_TWIN.md:100-102

\begin{table}[ht]
  \centering
  \small
  \begin{tabular}{ll}
    \hline
    Quantity & Value \\
    \hline
    Test cycles, errors, flagged cycles & \RESULT{counts} \\
    Base error rate on the test half & \RESULT{rate} \\
    Flag rate on the test half & \RESULT{rate} \\
    H1 difference and one-sided lower bound & \RESULT{H1 interval} \\
    H1 grade & \RESULT{pass, fail or inconclusive} \\
    Output-invariant test cycles & \RESULT{count} \\
    H2 difference and one-sided lower bound & \RESULT{H2 interval} \\
    H2 grade & \RESULT{pass, fail or inconclusive} \\
    H3 & \RESULT{run only if H1 passes} \\
    \hline
  \end{tabular}
  \caption{Registered outcomes of the I-EIP study. Every entry awaits the graded record.}
  \label{tab:ieip-results}
\end{table}

\RESULT{per-transform and per-layer error distributions}

\RESULT{flagged cycles with their scenarios, and the null transform's score per cycle}

\RESULT{named failure modes group\_symmetry\_break and layerwise\_drift, where they apply}

\section{What a pass or a failure would mean}

A pass on H1 would show that the internal check carries information about the classifier's
errors on this robot's inputs.
A pass on H2 would show that it carries information an output check does not.
Neither would change what is proved in Chapter~\ref{ch:lean}.
The containment does not depend on the classifier being right.
A gate built on the monitor would act only by withholding classified events, which can make the
robot ask more often but cannot grant it anything.
That is why H3 measures false clears and over-restriction together.

A failure would agree with the stated prior.
It would add another negative result for this family of monitors, on a different model and
task from the sealed families of CR-IEIP.

\RESULT{interpretation after grading}
